\bichapter{介绍和使用}{Introduction}\label{chap:intro}
\linespread{1.08}\selectfont

欢迎使用北京工业大学硕士研究生学位\LaTeX 论文模版（学术型硕士）。制作本模版的思路是因为Office for MacOS的体验极差，加
之本人又对word十分不熟练。网上目前没有现成的模版给我白嫖，本着开源精神和拙劣的\LaTeX 理解，参考部分模版（见致谢）赶工
出该模版的1.0版本。希望借助此模版，使用者能将本就在毕业季不多的精力专注于文字而不是行距、行高、段落等繁琐的格式上。接下来将介绍模版的使用方法。

本模版中大量格式参考《北京工业大学研究生学位论文撰写规范》和北京工业大学博士论文\LaTeX 模版（见致谢）。但仍有小部分
细节格式未确认（见readme todo list）。关于格式定制问题，欢迎issue 至此模版的git仓库。

一些需要你手动设置的内容：

小节之前插入下行语句保证间隔
\begin{center}
  \begin{minipage}{0.5\textwidth}
    \begin{Verbatim}[frame=single]
    \hspace*{\fill}
    \end{Verbatim}
  \end{minipage}
\end{center}

\$M Word的行距计算涉及到不同字体类型，多倍行距单倍行距等内容，甚至可能在不同操作系统下不一致，这和\LaTeX 的基线算法十
分有出入，因此，为了和让人很迷糊的\$M Word 看起来尽可能视觉效果一致，需要手动在每个正文文件前设置如下内容。

\begin{center}
  \begin{minipage}{0.5\textwidth}
    \begin{Verbatim}[frame=single]
\linespread{1.08}\selectfont
    \end{Verbatim}
  \end{minipage}
\end{center}


\hspace*{\fill}

\bisection{如何使用}{howtouse}\label{sec:fst}
\subsection{模版文件总览}
首先要明确的是，\LaTeX 出于安全的考虑，在unix环境下，往上一级目录写入编译时文件（如aux文件）时会被禁止。为了最大兼容
性，当前版本的模版将所有的文件都放置在同一级目录下。（不排除某些tex的特殊编译设置或单独设置的环境变量）

大多数用户使用时，都不必关注所有文件，只需了解如下几个文件：

\begin{center}
  \begin{minipage}{0.95\textwidth}
    \begin{Verbatim}[frame=single]
    主文档 : main.tex         中文摘要 : cabstract.tex
    英文摘要 : eabstract.tex  章节文档 : chapt*.tex
    引用文档 : reference.bib  用户出版物文档 : publication.tex
    致谢文档 : acknowledge.tex
    \end{Verbatim}
  \end{minipage}
\end{center}

以上文档使用时，基本上只需填入文字内容即可，不需要考虑格式。

% 文字\cite{Boutsidis2011}
\subsubsection{主文档}

主文档主要用于记录用户基本信息，组织论文结构。在主文档中，如下内容需要用户自行填写或更改：

\begin{center}
  \begin{minipage}{0.65\textwidth}
    \begin{Verbatim}[frame=single]
      中文图书分类号 clc
      UDC分类号 udc
      学校代码 schoolcode
      密级 secretlevel
      论文中文题目 ctitle
      作者姓名 cauthor
      学科 cdiscipline
      指导教师 csupervisior
      指导教师职称 csupervstitle
      论文提交/答辩日期 cdate
      学号 sutdentid
      论文英文题目 etitle
      研究方向 cmajor
      申请学位 cdegree
      所在单位 ccollege
      授予学位单位 corganization
    \end{Verbatim}
  \end{minipage}
\end{center}

主文档正文部分组织了论文结构，论文正文内容由\verb |chapt1.tex|, \verb|chapt2.tex|等组织。用户自行添加，参考章节文档一节。%（todo）

\subsubsection{中文摘要}

中文摘要内容由\verb|cabstract|环境包裹，两个段落之间以大于1行的空行作为分割。

在中文文档中，有\verb|ckeywords|环境包裹的关键词，以"；"分割。

\subsubsection{英文摘要}

 英文摘要与中文摘要一致，由\verb|eabstract|环境包裹。

 类似的，也有\verb|ekeywords|环境包裹的关键词，但以","分割。

 \subsubsection{章节文档}
 章节文档是论文的主要内容，在章节文档中，包含不同的小结，本模版共4级小结，分别由\verb|bisection|, \verb|subsection|, \verb|subsubsection|,  \verb|subsubsubsection| 标记。
 其中，目录会记录至第二级小结，即\verb|subsection|。具体书写格式参考本文档源码（\verb|chapt1.tex|）。

\subsubsection{引用文档}

引用文档格式为正常\verb|bib file|。

\subsubsection{出版物文档}

直接在\verb|publication|环境中添加即可。

\subsubsection{致谢文档}

同引用文档。

\bisection{基本语法}
下面介绍一些基本语法。可自行修改。
\subsection{如何插入图片并引用}

\begin{center}
  \begin{minipage}{0.95\textwidth}
    \begin{Verbatim}[frame=single]
\begin{figure}[htpb]
  \centering
  \includegraphics[width=0.8\textwidth]{bjut_logo_color.pdf}
  \bicaption{中文图题}{English Figure Title}
  \label{fig:logo}
\end{figure}

% 使用\cref{<label>}引用: 如\cref{fig:logo}所示
    \end{Verbatim}
  \end{minipage}
\end{center}

\begin{figure}[htpb]
  \centering
  \includegraphics[width=0.8\textwidth]{bjut_logo_color.pdf}
  \bicaption{中文图题}{English Figure Title}
  \label{fig:logo}
\end{figure}
“如\cref{fig:logo}所示”

\subsection{如何插入表格并引用}

\begin{center}
  \begin{minipage}{0.95\textwidth}
    \small
    \begin{Verbatim}[frame=single]
        \begin{table}[htpb]
          \centering
          \bicaption[]{三线表}{A Table Example}
          \label{tab:1st}
          \begin{tabular}{@{}llr@{}}
            \toprule
            \multicolumn{2}{c}{Item} & \\
            \cmidrule(r){1-2}
            Animal    & Description  & Price(\$) \\
            \midrule
            Gnat      & pergram      & 13.65 \\
            & each         & 0.01 \\
            Gnu       & stuffed      & 92.50 \\
            Emu       & stuffed      & 33.33 \\
            Armadillo & frozen       & 8.99 \\
            \bottomrule
          \end{tabular}
        \end{table}
        % 引用方法与图片一致 如\cref{tab:1st}所示
    \end{Verbatim}
  \end{minipage}
\end{center}

        \begin{table}[htpb]
          \centering
          \bicaption[]{三线表}{A Table Example}
          \label{tab:1st}
          \begin{tabular}{@{}llr@{}}
            \toprule
            \multicolumn{2}{c}{Item} & \\
            \cmidrule(r){1-2}
            Animal    & Description  & Price(\$) \\
            \midrule
            Gnat      & pergram      & 13.65 \\
            & each         & 0.01 \\
            Gnu       & stuffed      & 92.50 \\
            Emu       & stuffed      & 33.33 \\
            Armadillo & frozen       & 8.99 \\
            \bottomrule
          \end{tabular}
        \end{table}

“如\cref{tab:1st}所示”


\subsection{如何插入公式}

注意，插入公式时，公式环境和前后文字之间不应该有空行。

\begin{center}
  \begin{minipage}{0.85\textwidth}
    \begin{Verbatim}[frame=single]
      \begin{equation}\label{eq:eq1}
        e^{\pi i}+1 = 0
      \end{equation}
    \end{Verbatim}
  \end{minipage}
\end{center}

      \begin{equation}\label{eq:eq1}
        e^{\pi i}+1 = 0
      \end{equation}

\subsection{如何插入算法}

\begin{center}
  \begin{minipage}{0.85\textwidth}
    \begin{Verbatim}[frame=single]
\begin{algorithm}[htbp]
  \caption{1+1等于几}
  \label{alg:algorithm}
    \begin{algorithmic}[1]
      \REQUIRE $1$.
      \ENSURE $1$.
      \STATE $i \leftarrow 1$
      \FOR {$i \leq 1$}
        \STATE $i \leftarrow i + 1$;
      \ENDFOR
      \STATE Return $i$.
    \end{algorithmic}
\end{algorithm}
    \end{Verbatim}
  \end{minipage}
\end{center}

\begin{algorithm}[htbp]
\caption{1+1等于几}
\label{alg:algorithm}
\begin{algorithmic}[1]
\REQUIRE $1$.
\ENSURE $1$.
\STATE $i \leftarrow 1$
\FOR {$i \leq 1$}
\STATE $i \leftarrow i + 1$;
\ENDFOR
\STATE Return $i$.
\end{algorithmic}
\end{algorithm}
